\scanpage{211}

on trouve que

\[
P^{cip\,\unsure{0|o}}_{X/S}\xleftarrow{\sim}P^{cip\,0}_{A^1/S}\simeq P^{cip\,0}_{A^0/S}\simeq B^{(p/S)}\tag*{(3.7)}
\]

i.e. on trouve la m\textsuperscript{e} chose que \unsure{sur|dans} le cas d'un schéma abélien. \inserted{\uline{Remarque}} D'ailleurs, faisant \deleted{\unsure{le}} maintnt $X=A$, on sait que $B^{\unsure{(p/S)|cp0}}$, regardé comme extension de $B$ par un grpe \deleted{$\mathrm{Ker}\ V$}

\[
\alpha_B=\mathrm{Ker}(V_B:B^{p/S}\to B),
\]

qui est un groupe affine sur $S$ annulé par $V$, est l'extension \uline{universelle} de cette nature [du moins \deleted{\unsure{sur un}} si $S$ est le spectre d'un corps parfait].

$D(\alpha_B)=\mathrm{Ker}\,F_{B^*}$ \inserted{\deleted{\unsure{donc}} $F_A$} est un groupe infinitésimal de hauteur 1, dont l'algèbre de Lie est $\underline{\mathrm{Lie}}(B^*)$, et on sait que $\underline{\mathrm{Lie}}(B^*)^\vee$ est l'\uline{enveloppe vectorielle} de ce groupe fini; donc on voit que l'extension vectorielle universelle est déduite par $\alpha_B\longrightarrow\text{enveloppe vectorielle}\simeq\underline{\mathrm{Lie}}(B^*)^\vee$, de l'extension universelle par des groupes annulés par $V$.

4. Supposons que $p$ soit nilpotent sur $S$, sans plus. Soit $P^{cin}\subset P^{ci}$, défini par les connexions à courbure \deleted{nulle}\uline{nilpotentes}. Si $X/S$ est lisse, \deleted{et $S$ réduit, alors il s'agit de connexions} \inserted{les valeurs des foncteurs préc.\ sur les $S'$ \uline{réduits}} sont les \deleted{\unsure{mêmes}} que \deleted{\unsure{les}} ceux de $P^{cip}$; \deleted{et pour $S'$} quelcon-

% archive sheet number: 210
